\documentclass{article}
\usepackage[utf8]{inputenc}    
\usepackage[T1]{fontenc}       
\usepackage{lmodern}           
\usepackage{hyperref}
\hypersetup{
    colorlinks=true,       % false: boxed links; true: colored links
    linkcolor=blue,        % color of internal links (sections, etc.)
    urlcolor=cyan,         % color of external links
    linktoc=all            % 'all' makes both text and page number clickable
}
\usepackage{amsmath}   
\usepackage{amssymb}   
\usepackage{geometry}  
\usepackage{enumerate} 
\usepackage{xcolor}  
\usepackage{amsthm}
\usepackage{pdfpages}
\usepackage{mathrsfs}
\newtheorem{theorem}{Theorem}[section]
% \newtheorem{lemma}[theorem]{Lemma}
\newtheorem{corollary}[theorem]{Corollary}
% \newtheorem{definition}[theorem]{Definition}
\newtheorem{example}[theorem]{Example}
\usepackage{listings}  
\usepackage{tikz-cd}
\usetikzlibrary{positioning}
\usepackage{forest}     
\usepackage{xcolor}
\usepackage{tcolorbox}
\tcbuselibrary{theorems}

% remove/comment out the old line:
% \newtheorem{definition}[theorem]{Definition}

\usepackage{xcolor}
\usepackage{tcolorbox}
\tcbuselibrary{breakable}
\newtheorem{definition}[theorem]{Definition}
\let\olddefinition\definition
\let\endolddefinition\enddefinition
\renewenvironment{definition}
  {\begin{tcolorbox}[
      colback=red!10,
      colframe=red!50!black,
      boxrule=0.6pt,
      arc=2pt,
      left=6pt,right=6pt,top=4pt,bottom=4pt,
      breakable
    ]\olddefinition}
  {\endolddefinition\end{tcolorbox}}

\newtheorem{lemma}[theorem]{Lemma}
\let\oldlemma\lemma
\let\endoldlemma\endlemma
\renewenvironment{lemma}
  {\begin{tcolorbox}[
      colback=green!8,
      colframe=green!50!black,
      boxrule=0.6pt,
      arc=2pt,
      left=6pt,right=6pt,top=4pt,bottom=4pt
    ]\oldlemma}
  {\endoldlemma\end{tcolorbox}}

  \let\oldtheorem\theorem
\let\endoldtheorem\endtheorem
\renewenvironment{theorem}
  {\begin{tcolorbox}[
      colback=yellow!6,
      colframe=yellow!60!black,
      boxrule=0.6pt,
      arc=2pt,
      left=6pt,right=6pt,top=4pt,bottom=4pt,
      parbox=false
    ]\oldtheorem}
  {\endoldtheorem\end{tcolorbox}}
\usepackage{bussproofs}
\usetikzlibrary{arrows.meta}
\usetikzlibrary{arrows.meta,decorations.pathreplacing,calc}
\lstset{frame=tb,
  language=C,
  aboveskip=3mm,
  belowskip=3mm,
  showstringspaces=false,   
  columns=flexible,
  basicstyle={\small\ttfamily},
  numbers=none,
  numberstyle=\tiny\color{gray},
  keywordstyle=\color{blue},
  commentstyle=\color{brown},
  stringstyle=\color{orange},
  breaklines=true,
  breakatwhitespace=true,
  tabsize=3
}
\geometry{top=0.5in, bottom=0.5in, left=1in, right=1in}
% ctrl shift D to toggle night 
\begin{document}

\title{HW4}
\author{Wang Xiyu A0282500R}
\date{}
\maketitle

\section*{Question 1}

\begin{lemma}[Relevance Lemma for FOL]
Let $\phi$ be a FOL formula, let $\mathfrak{A}$ be a structure and let $\alpha_1,\alpha_2$ be
variable assignments such that for all $x \in FVar(\phi)$,
\[
\alpha_1(x) = \alpha_2(x).
\]
Then
\[
\mathfrak{A},\alpha_1 \models \phi \iff \mathfrak{A},\alpha_2 \models \phi.
\]
\end{lemma}

\textbf{Question.} Prove Lemma 1.
\begin{proof}
    By symmetry, we only need to prove if for all $x\in FVar(\phi), \alpha_1(x) = \alpha_2(x)$
    then
    \[
\mathfrak{A},\alpha_1 \models \phi \implies \mathfrak{A},\alpha_2 \models \phi.
\]
\begin{lemma} If
    \[\forall x \in Var(t), \alpha_1(x) = \alpha_2(x)\]
    then $Val_{\mathfrak{A}, \alpha_1}(t) = Val_{\mathfrak{A}, \alpha_2}(t)$
\end{lemma}
\subsubsection*{Base case: }
\begin{enumerate}
    \item $t \equiv c$: $I(c) = I(c) \implies Val_{\mathfrak{A}, \alpha_1}(c) = Val_{\mathfrak{A}, \alpha_2}(c)$
    \item $t \equiv x$: By definition, $Val_{\mathfrak{A}, \alpha_1}(x) = \alpha_1(x) = \alpha_2(x) = Val_{\mathfrak{A}, \alpha_2}(x)$
\end{enumerate}
\subsubsection*{Inductive step}
Let
\[
t \equiv f(t_1,\dots,t_n).
\]
\begin{align*}
Val_{\mathfrak{A},\alpha_1}(t)
&= I_{\mathfrak{A}}(f)\bigl(Val_{\mathfrak{A},\alpha_1}(t_1),\dots,Val_{\mathfrak{A},\alpha_1}(t_n)\bigr).
\end{align*}
By IH, for all $i$
\[
Val_{\mathfrak{A},\alpha_1}(t_i)=Val_{\mathfrak{A},\alpha_2}(t_i)
\]
\[
\bigl(Val_{\mathfrak{A},\alpha_1}(t_1),\dots,Val_{\mathfrak{A},\alpha_1}(t_n)\bigr)
=
\bigl(Val_{\mathfrak{A},\alpha_2}(t_1),\dots,Val_{\mathfrak{A},\alpha_2}(t_n)\bigr).
\]
\begin{align*}
Val_{\mathfrak{A},\alpha_1}(t)
&= I_{\mathfrak{A}}(f)\bigl(Val_{\mathfrak{A},\alpha_1}(t_1),\dots\bigr) \\
&= I_{\mathfrak{A}}(f)\bigl(Val_{\mathfrak{A},\alpha_2}(t_1),\dots\bigr) \\
&= Val_{\mathfrak{A},\alpha_2}(t).
\end{align*}
Now we can use the lemma to prove the previosu lemma: 
\subsubsection*{Base cases}

\begin{enumerate}
    \item $\phi \equiv t_1 = t_2$.
    Suppose
    \[
    \mathfrak{A},\alpha_1 \models t_1=t_2.
    \]
    Then we have
    \[
    Val_{\mathfrak{A},\alpha_1}(t_1)=Val_{\mathfrak{A},\alpha_1}(t_2).
    \]
    \[
    FVar(t_1=t_2)=Var(t_1)\cup Var(t_2),
    \]
    \[
    Val_{\mathfrak{A},\alpha_1}(t_1)=Val_{\mathfrak{A},\alpha_2}(t_1)
    \quad\text{and}\quad
    Val_{\mathfrak{A},\alpha_1}(t_2)=Val_{\mathfrak{A},\alpha_2}(t_2).
    \]
    \[
    Val_{\mathfrak{A},\alpha_2}(t_1)=Val_{\mathfrak{A},\alpha_2}(t_2) \implies 
    \mathfrak{A},\alpha_2 \models t_1=t_2.
    \]

    \item $\phi \equiv R(t_1,\dots,t_k)$.

    Suppose
    \[
    \mathfrak{A},\alpha_1 \models R(t_1,\dots,t_k).
    \]
    Then
    \[
    \bigl(Val_{\mathfrak{A},\alpha_1}(t_1),\dots,Val_{\mathfrak{A},\alpha_1}(t_k)\bigr)\in I(R).
    \]
    \[
    FVar(R(t_1,\dots,t_k))=Var(t_1)\cup \cdots \cup Var(t_k),
    \]
    by the term lemma, 
    \[
    \forall i, Val_{\mathfrak{A},\alpha_1}(t_i)=Val_{\mathfrak{A},\alpha_2}(t_i).
    \]
    \[
    \bigl(Val_{\mathfrak{A},\alpha_1}(t_1),\dots,Val_{\mathfrak{A},\alpha_1}(t_k)\bigr) = \bigl(Val_{\mathfrak{A},\alpha_2}(t_1),\dots,Val_{\mathfrak{A},\alpha_2}(t_k)\bigr) \in I(R)
    \]
    \[
    \mathfrak{A},\alpha_2 \models R(t_1,\dots,t_k).
    \]
\end{enumerate}

\subsubsection*{Inductive cases}

\begin{enumerate}
    \item $\phi \equiv \neg \psi$.

    Suppose
    \[
    \mathfrak{A},\alpha_1 \models \neg \psi.
    \]
    Then
    \[
    \mathfrak{A},\alpha_1 \not\models \psi.
    \]
    \[
    FVar(\neg\psi)=FVar(\psi),
    \]
    \[
    \mathfrak{A},\alpha_2 \models \psi \implies \mathfrak{A},\alpha_1 \models \psi.
    \]
    \[
    \mathfrak{A},\alpha_2 \not\models \psi,
    \]
    \[
    \mathfrak{A},\alpha_2 \models \neg \psi.
    \]

    \item $\phi \equiv \psi_1 \lor \psi_2$.

    \[
    \mathfrak{A},\alpha_1 \models \psi_1 \lor \psi_2.
    \]
    then either
    \[
    \mathfrak{A},\alpha_1 \models \psi_1
    \quad\text{or}\quad
    \mathfrak{A},\alpha_1 \models \psi_2.
    \]
    Also,
    \[
    FVar(\psi_1\lor\psi_2)=FVar(\psi_1)\cup FVar(\psi_2).
    \]
    by IH
    \[
    \mathfrak{A},\alpha_2 \models \psi_1
    \quad\text{or}\quad
    \mathfrak{A},\alpha_2 \models \psi_2.
    \]
    Therefore,
    \[
    \mathfrak{A},\alpha_2 \models \psi_1 \lor \psi_2.
    \]

    \item $\phi \equiv \exists x\,\psi$.

    Suppose
    \[
    \mathfrak{A},\alpha_1 \models \exists x\,\psi.
    \]
    Then there exists some $u\in U$ such that
    \[
    \mathfrak{A},\alpha_1[x\mapsto u]\models \psi.
    \]
    w can show that $\alpha_1[x\mapsto u]$ and $\alpha_2[x\mapsto u]$ agree on all variables in $FVar(\psi)$.

    fix a $y\in FVar(\psi)$,

    \begin{itemize}
        \item If $y=x$, then
        \[
        \alpha_1[x\mapsto u](y)=u=\alpha_2[x\mapsto u](y).
        \]

        \item If $y\neq x$, then $y \in FVar(\psi)\setminus\{x\} = FVar(\exists x, \psi)$
        \[
        \alpha_1[x\mapsto u](y)=\alpha_1(y) = \alpha_2(y)=\alpha_2[x\mapsto u](y).
        \]
    \end{itemize}

    Thus, $\forall y\in FVar(\psi)$,
    \[
    \alpha_1[x\mapsto u](y)=\alpha_2[x\mapsto u](y).
    \]
    By IH
    \[
    \mathfrak{A},\alpha_2[x\mapsto u]\models \psi.
    \]
    Hence
    \[
    \mathfrak{A},\alpha_2 \models \exists x\,\psi.
    \]

    \item $\phi \equiv \forall x\,\psi$.

    Suppose
    \[
    \mathfrak{A},\alpha_1 \models \forall x\,\psi.
    \]
    then for every $u\in U$,
    \[
    \mathfrak{A},\alpha_1[x\mapsto u]\models \psi.
    \]
    Fix a $u\in U$. We can show that $\alpha_1[x\mapsto u]$ and $\alpha_2[x\mapsto u]$ agree on all variables in $FVar(\psi)$.

    fix a $y\in FVar(\psi)$.

    \begin{itemize}
        \item If $y=x$, then
        \[
        \alpha_1[x\mapsto u](y)=u=\alpha_2[x\mapsto u](y).
        \]

        \item If $y\neq x$, then $y \in FVar(\psi)\setminus\{x\} = FVar(\exists x, \psi)$
        \[
        \alpha_1[x\mapsto u](y)=\alpha_1(y) = \alpha_2(y)=\alpha_2[x\mapsto u](y).
        \]
    \end{itemize}
    by IH,
    \[
    \mathfrak{A},\alpha_2[x\mapsto u]\models \psi.
    \]
    Since $u\in U$ is arbitrary,
    \[
    \mathfrak{A},\alpha_2 \models \forall x\,\psi.
    \]
\end{enumerate}
\end{proof}

\section*{Question 2}

\begin{definition}[Ground Terms, Atoms and Formulae]
A term over a signature $\Sigma$ is said to be a \emph{ground term} if it does not contain any
occurrence of a variable. The set of all ground terms over $\Sigma$ is denoted by
\[
\mathrm{Terms}_\Sigma .
\]

A formula is said to be \emph{ground} if it contains no variables (and therefore no quantifiers).

A formula is said to be a \emph{ground atom} if it is both ground and atomic.

We denote by
\[
\mathrm{Atoms}_\Sigma
\]
the set of all ground atoms, and by
\[
\mathrm{FORM}_\Sigma
\]
the set of all ground formulae over $\Sigma$.
\end{definition}

\begin{definition}[Ground Instantiation]
Let
\[
\phi \equiv \forall x_1 \forall x_2 \cdots \forall x_n \psi
\]
be a universally quantified FOL sentence in prenex normal form.

A \emph{ground instantiation} of $\phi$ is obtained by replacing each occurrence of $x_i$
in $\psi$ by a ground term $t_i$.

Formally, there exists a mapping
\[
T : V \rightarrow \mathrm{Terms}_\Sigma
\]
such that
\[
\varphi \equiv \psi[x_1 \rightarrow T(x_1)]\cdots[x_n \rightarrow T(x_n)].
\]

The set of all ground instantiations of $\phi$ is denoted
\[
\mathrm{Ground}(\phi).
\]

For a set $\Gamma$ of universally quantified sentences,
\[
\mathrm{Ground}(\Gamma) =
\bigcup_{\phi \in \Gamma} \mathrm{Ground}(\phi).
\]
\end{definition}

\begin{definition}[Propositional Satisfiability]
A \emph{ground atomic valuation} over $\Sigma$ is a mapping
\[
\tau : \mathrm{Atoms}_\Sigma \rightarrow \{\mathrm{true},\mathrm{false}\}.
\]

This valuation is extended to ground formulae using the relation $\models_{Gr}$ defined by:

\[
\tau \models_{Gr} \psi
\quad \text{iff} \quad
\tau(\psi) = \mathrm{true}
\qquad \text{if } \psi \in \mathrm{Atoms}_\Sigma
\]

\[
\tau \models_{Gr} \neg \phi
\quad \text{iff} \quad
\tau \not\models_{Gr} \phi
\]

\[
\tau \models_{Gr} (\phi_1 \vee \phi_2)
\quad \text{iff} \quad
\tau \models_{Gr} \phi_1
\text{ or }
\tau \models_{Gr} \phi_2 .
\]

A set of ground FOL formulae $\Gamma$ is \emph{propositionally satisfiable} if there exists
a ground atomic valuation $\tau$ such that
\[
\tau \models_{Gr} \phi
\quad \text{for every } \phi \in \Gamma .
\]
\end{definition}

\begin{lemma}
If $\Gamma$ is satisfiable then $\mathrm{Ground}(\Gamma)$ is propositionally satisfiable.
\end{lemma}

\textbf{Question.} Prove Lemma 2.
\begin{proof}
    $\Gamma$ is satisfiable, there exists a structure $\mathfrak A$ such that
    \[
    \mathfrak A \models \Gamma.
    \]
    Define ground atom valuation
    \[
    \tau_{\mathfrak A} : Atoms_\Sigma \to \{\mathrm{true},\mathrm{false}\}
    \]
    \[
    \tau_{\mathfrak A}(a)=\mathrm{true}
    \iff
    \mathfrak A \models a
    \]
    for every ground atom $a$ over $\Sigma$.
    Now we show that for every ground formula $\phi$,
    \[
    \mathfrak A \models \phi \implies \tau_{\mathfrak A}\models_{Gr}\phi.
    \]
    Induction on the structure of $\phi$    
    \subsubsection*{Base case: $\chi$ is a ground atom}
    trivial    
    \subsubsection*{Inductive cases}
    \begin{enumerate}
        \item $\psi \equiv \neg \phi$.
        Suppose
        \[
        \mathfrak A \models \neg \phi \implies
        \mathfrak A \not\models \phi
        \]
        By IH,
        \[
        \mathfrak A \models \phi \implies \tau_{\mathfrak A}\models_{Gr}\phi
        \]
        \[
        \tau_{\mathfrak A}\not\models_{Gr}\phi
        \]
        \[
        \tau_{\mathfrak A}\models_{Gr}\neg\phi
        \]
    
        \item $\chi \equiv \phi_1 \lor \phi_2$.
    
        Suppose
        \[
        \mathfrak A \models \phi_1 \lor \phi_2
        \]
        Then 
        \[
        \mathfrak A \models \phi_1
        \quad\text{or}\quad
        \mathfrak A \models \phi_2
        \]
        by IH
        \[
        \tau_{\mathfrak A}\models_{Gr}\phi_1
        \quad\text{or}\quad
        \tau_{\mathfrak A}\models_{Gr}\phi_2
        \]
        \[
        \tau_{\mathfrak A}\models_{Gr}\phi_1 \lor \phi_2
        \]
    \end{enumerate}
    Fix any $\chi \in Ground(\Gamma)$. Then there exists
    \[
    \varphi=\forall x_1\cdots\forall x_n\,\psi \in \Gamma
    \]
    such that
    \[
    \chi=\psi[x_1\mapsto t_1]\cdots[x_n\mapsto t_n]
    \]
    for ground terms $t_1,\dots,t_n$\\
    $\mathfrak A \models \Gamma \implies \mathfrak A \models \varphi$\\
    For $1 \leq i \leq n$
    \[
    u_i = Val_{\mathfrak A}(t_i)
    \]
    \[
    \mathfrak A,\alpha[x_1\mapsto u_1]\cdots[x_n\mapsto u_n]\models \psi.
    \]
    By substitution lemma, which we will prove later
    \[
    \mathfrak A \models \psi[x_1\mapsto t_1]\cdots[x_n\mapsto t_n].
    \]
    \[
    \mathfrak A \models \chi.
    \]
    \[
    \tau_{\mathfrak A}\models_{Gr}\chi.
    \]
    
    Since $\chi\in Ground(\Gamma)$ is arbitrarily chosen, we have
    \[
    \forall \chi \in Ground(\Gamma), \tau_{\mathfrak{A}} \models_{Gr} \chi
    \]
    \[
    \tau_{\mathfrak A}\models_{Gr} Ground(\Gamma).
    \]
    Therefore $Ground(\Gamma)$ is propositionally sat.
    \end{proof}
\begin{lemma}[Substitution lemma]
    \[\mathfrak{A}\models \psi[x \mapsto t] \iff \mathfrak{A}, \alpha[x \mapsto Val_{\mathfrak{A}(t)}] \models \psi\]
\end{lemma}
\begin{lemma}
    \[Val_{\mathfrak{A}, \alpha}(t[x\mapsto t']) = Val_{\mathfrak{A}, \alpha[x \mapsto Val_{\mathfrak{A}, \alpha}(t')]}(t)\]
\end{lemma}
\begin{proof}
\subsubsection*{Base case: }
\begin{enumerate}
    \item $t = c \in C$: 
    \[Val_{\mathfrak{A}, \alpha}(t[x \mapsto t']) = Val_{\mathfrak{A}, \alpha}(c) = I(c)\]
    \[Val_{\mathfrak{A}, \alpha[x \mapsto Val_{\mathfrak{A}, \alpha}(t')]}(t) = Val_{\mathfrak{A}, \alpha[x \mapsto Val_{\mathfrak{A}, \alpha}(t')]}(c) = I(c)\]
    \item $t = y \in V$:\\
    Case 1: $x = y$
    \[
    \implies y[x \mapsto t'] = t'
    \]
    \[
    Val_{\mathfrak{A}, \alpha}(y[x \mapsto t']) = Val_{\mathfrak{A}, \alpha}(t')
    \]
    \[
    Val_{\mathfrak{A}, \alpha[x \mapsto Val_{\mathfrak{A},\alpha}(t')]}(y)
    = \alpha[x \mapsto Val_{\mathfrak{A},\alpha}(t')](y)
    = Val_{\mathfrak{A},\alpha}(t')
    \]
    Case 2: $x \neq y$
    \[
    \implies y[x \mapsto t'] = y
    \]
    \[
    Val_{\mathfrak{A}, \alpha}(y[x \mapsto t']) = Val_{\mathfrak{A}, \alpha}(y) = \alpha(y)
    \]
    \[
    Val_{\mathfrak{A}, \alpha[x \mapsto Val_{\mathfrak{A},\alpha}(t')]}(y)
    = \alpha[x \mapsto Val_{\mathfrak{A},\alpha}(t')](y)
    = \alpha(y)
    \]
\end{enumerate}
\subsubsection*{Inductive step: $t = f(t_1, \dots)$}
\[
    t[x \mapsto t'] = 
    f(t_1[x \mapsto t'], t_2[x \mapsto t'], \dots)
\]
\[
    Val_{\mathfrak{A}, \alpha}(t[x \mapsto t']) = 
    Val_{\mathfrak{A}, \alpha}(f(t_1[x \mapsto t'], t_2[x \mapsto t'], \dots)) =
\]
\[
    I(f)(Val_{\mathfrak{A}, \alpha}(t_1[x \mapsto t']), Val_{\mathfrak{A}, \alpha}(t_2[x \mapsto t']), \dots)
\]
The other side
\[
    Val_{\mathfrak{A}, \alpha[x \mapsto Val_{\mathfrak{A},\alpha}(t')]}(f(t_1, t_2, \dots)) = 
    I(f)(Val_{\mathfrak{A}, \alpha}(t_1[x \mapsto t']), Val_{\mathfrak{A}, \alpha}(t_2[x \mapsto t']), \dots)
\]
\end{proof}
\begin{proof}
    We prove the substitution lemma for formulae by structural induction on $\psi$.
    
    \subsubsection*{Base case: $\psi \equiv t_1 = t_2$}
    \[
    \psi[x \mapsto t] \equiv t_1[x \mapsto t] = t_2[x \mapsto t]
    \]
    Thus
    \[
    \mathfrak{A}, \alpha \models \psi[x \mapsto t]
    \iff
    Val_{\mathfrak{A}, \alpha}(t_1[x \mapsto t])
    =
    Val_{\mathfrak{A}, \alpha}(t_2[x \mapsto t]).
    \]
    By the substitution lemma for terms,
    \[
    Val_{\mathfrak{A}, \alpha}(t_1[x \mapsto t])
    =
    Val_{\mathfrak{A}, \alpha[x \mapsto Val_{\mathfrak{A}, \alpha}(t)]}(t_1)
    \]
    and
    \[
    Val_{\mathfrak{A}, \alpha}(t_2[x \mapsto t])
    =
    Val_{\mathfrak{A}, \alpha[x \mapsto Val_{\mathfrak{A}, \alpha}(t)]}(t_2).
    \]
    Hence
    \[
    Val_{\mathfrak{A}, \alpha[x \mapsto Val_{\mathfrak{A}, \alpha}(t)]}(t_1)
    =
    Val_{\mathfrak{A}, \alpha[x \mapsto Val_{\mathfrak{A}, \alpha}(t)]}(t_2),
    \]
    that is,
    \[
    \mathfrak{A}, \alpha[x \mapsto Val_{\mathfrak{A}, \alpha}(t)] \models t_1 = t_2.
    \]
    
    \subsubsection*{Inductive steps}
    \begin{enumerate}
        \item $\psi \equiv R(t_1, \dots, t_n)$.
    
        \[
        \psi[x \mapsto t] \equiv R(t_1[x \mapsto t], \dots, t_n[x \mapsto t]).
        \]
        Then
        \[
        \mathfrak{A}, \alpha \models \psi[x \mapsto t]
        \iff
        \bigl(Val_{\mathfrak{A}, \alpha}(t_1[x \mapsto t]), \dots, Val_{\mathfrak{A}, \alpha}(t_n[x \mapsto t])\bigr) \in I(R).
        \]
        By the substitution lemma for terms, for each $i$,
        \[
        Val_{\mathfrak{A}, \alpha}(t_i[x \mapsto t])
        =
        Val_{\mathfrak{A}, \alpha[x \mapsto Val_{\mathfrak{A}, \alpha}(t)]}(t_i).
        \]
        \[
        \bigl(Val_{\mathfrak{A}, \alpha[x \mapsto Val_{\mathfrak{A}, \alpha}(t)]}(t_1), \dots,
        Val_{\mathfrak{A}, \alpha[x \mapsto Val_{\mathfrak{A}, \alpha}(t)]}(t_n)\bigr) \in I(R),
        \]
        \[
        \mathfrak{A}, \alpha[x \mapsto Val_{\mathfrak{A}, \alpha}(t)] \models R(t_1, \dots, t_n).
        \]
    
        \item $\psi \equiv \neg \phi$.
    
        \[
        \mathfrak{A}, \alpha \models (\neg \phi)[x \mapsto t]
        \iff
        \mathfrak{A}, \alpha \models \neg (\phi[x \mapsto t])
        \]
        \[
        \iff
        \mathfrak{A}, \alpha \not\models \phi[x \mapsto t].
        \]
        By IH,
        \[
        \mathfrak{A}, \alpha \models \phi[x \mapsto t]
        \iff
        \mathfrak{A}, \alpha[x \mapsto Val_{\mathfrak{A}, \alpha}(t)] \models \phi.
        \]
        \[
        \mathfrak{A}, \alpha[x \mapsto Val_{\mathfrak{A}, \alpha}(t)] \not\models \phi,
        \]
        \[
        \mathfrak{A}, \alpha[x \mapsto Val_{\mathfrak{A}, \alpha}(t)] \models \neg \phi.
        \]
    
        \item $\psi \equiv \phi_1 \lor \phi_2$.
    
        \[
        \mathfrak{A}, \alpha \models (\phi_1 \lor \phi_2)[x \mapsto t]
        \iff
        \mathfrak{A}, \alpha \models \phi_1[x \mapsto t] \lor \phi_2[x \mapsto t].
        \]
        \[
        \mathfrak{A}, \alpha \models \phi_1[x \mapsto t]
        \quad \text{or} \quad
        \mathfrak{A}, \alpha \models \phi_2[x \mapsto t].
        \]
        By IH
        \[
        \mathfrak{A}, \alpha[x \mapsto Val_{\mathfrak{A}, \alpha}(t)] \models \phi_1
        \quad \text{or} \quad
        \mathfrak{A}, \alpha[x \mapsto Val_{\mathfrak{A}, \alpha}(t)] \models \phi_2.
        \]
        \[
        \mathfrak{A}, \alpha[x \mapsto Val_{\mathfrak{A}, \alpha}(t)] \models \phi_1 \lor \phi_2.
        \]
    \end{enumerate}
    \end{proof}
\section*{Question 3}

\begin{definition}[Herbrand Structure]
Let $\Sigma=(C,F,R)$ be a signature with $C\neq \emptyset$.

The $\Sigma$-Herbrand structure is
\[
\mathcal{A}_\Sigma=(U_\Sigma,I_\Sigma)
\]
defined as follows.

\begin{itemize}

\item The universe is
\[
U_\Sigma=\mathrm{Terms}_\Sigma
\]
(the set of all ground terms).

\item Constants are interpreted naturally:
\[
I_\Sigma(c)=c.
\]

\item For every function symbol $f$ of arity $r$,
\[
I_\Sigma(f)(t_1,\dots,t_r)=f(t_1,\dots,t_r).
\]

\item Relation symbols may be interpreted arbitrarily.

\end{itemize}

\end{definition}

\begin{lemma}
If $\mathrm{Ground}(\Gamma)$ is propositionally satisfiable then $\Gamma$ is satisfiable.
\end{lemma}

\textbf{Question.} Prove Lemma 3.
\begin{proof}
    Suppose $Ground(\Gamma)$ is propositionally sat, there is a ground atomic valuatio $\tau$
    \[
    \forall \phi \in Ground(\Gamma), \tau \models_{Gr} \phi 
    \]
    Construct a Herbrand structure as per definition. For $I(R)$, we define as such:
    \[I_\Sigma^\tau(R) = \{(t_1, t_2, \dots, t_{arity(R)})  \mid \tau\models_{Gr} R(t_1, t_2, \dots, t_{arity(R)}) \}\]
    Now we show that for every ground formula $\psi$,
    \[
    \mathfrak{A}^\tau_\Sigma \models \psi \iff \tau \models_{Gr} \psi.
    \]
    
    \begin{lemma}
    For every ground term $t$,
    \[
    Val_{\mathfrak{A}^\tau_\Sigma}(t)=t.
    \]
    \end{lemma}
    
    \begin{proof}
    By structural induction on $t$.
    If $t\equiv c$ for some constant $c$, then
    \[
    Val_{\mathfrak{A}^\tau_\Sigma}(c)=I^\tau_\Sigma(c)=c.
    \]
    If $t\equiv f(t_1,\dots,t_r)$, then by IH
    \[
    Val_{\mathfrak{A}^\tau_\Sigma}(t_i)=t_i
    \quad\text{for each }i.
    \]
    \[
    Val_{\mathfrak{A}^\tau_\Sigma}(f(t_1,\dots,t_r))
    =
    I^\tau_\Sigma(f)(Val_{\mathfrak{A}^\tau_\Sigma}(t_1),\dots,Val_{\mathfrak{A}^\tau_\Sigma}(t_r))
    =
    I^\tau_\Sigma(f)(t_1,\dots,t_r)
    =
    f(t_1,\dots,t_r).
    \]
    \end{proof}
    \subsubsection*{Base cases}
    \begin{enumerate}
        \item $\psi \equiv R(t_1,\dots,t_r)$.
    
        By definition,
        \[
        \mathfrak{A}^\tau_\Sigma \models R(t_1,\dots,t_r)
        \iff
        (Val_{\mathfrak{A}^\tau_\Sigma}(t_1),\dots,Val_{\mathfrak{A}^\tau_\Sigma}(t_r))\in I^\tau_\Sigma(R).
        \]
        \[
        (t_1,\dots,t_r)\in I^\tau_\Sigma(R).
        \]
        By definition of $I^\tau_\Sigma(R)$, 
        \[
        \tau \models_{Gr} R(t_1,\dots,t_r).
        \]
    \end{enumerate}
    
    \subsubsection*{Inductive cases}
    
    \begin{enumerate}
        \item $\psi \equiv \neg \phi$.
        \[
        \mathfrak{A}^\tau_\Sigma \models \neg\phi
        \iff
        \mathfrak{A}^\tau_\Sigma \not\models \phi.
        \]
        By IH
        \[
        \mathfrak{A}^\tau_\Sigma \not\models \phi
        \iff
        \tau \not\models_{Gr} \phi.
        \]
        \[
        \mathfrak{A}^\tau_\Sigma \models \neg\phi
        \iff
        \tau \models_{Gr} \neg\phi.
        \]
    
        \item $\psi \equiv \phi_1\lor\phi_2$.
        \[
        \mathfrak{A}^\tau_\Sigma \models \phi_1\lor\phi_2
        \iff
        \mathfrak{A}^\tau_\Sigma \models \phi_1
        \text{ or }
        \mathfrak{A}^\tau_\Sigma \models \phi_2.
        \]
        By IH
        \[
        \tau \models_{Gr} \phi_1
        \text{ or }
        \tau \models_{Gr} \phi_2.
        \]
        \[
        \mathfrak{A}^\tau_\Sigma \models \phi_1\lor\phi_2
        \iff
        \tau \models_{Gr} \phi_1\lor\phi_2.
        \]
    \end{enumerate}
    Now fix a frmula 
\[
\varphi=\forall x_1\cdots\forall x_n\,\psi \in \Gamma.
\]
We show that
\[
\mathfrak{A}^\tau_\Sigma \models \varphi
\]
fix arbitrary $u_1,\dots,u_n \in U_\Sigma = Terms_\Sigma$
each $u_i$ is a ground term
\[
\chi \equiv \psi[x_1\mapsto u_1]\cdots[x_n\mapsto u_n]
\]
is a ground instantiation of $\varphi$, thus 
\[
\chi \in Ground(\Gamma).
\]
as $\tau$ propositionally satisfies $Ground(\Gamma)$
\[
\tau \models_{Gr} \chi.
\]
\[
\mathfrak{A}^\tau_\Sigma \models \chi.
\]
\[
\mathfrak{A}^\tau_\Sigma \models \psi[x_1\mapsto u_1]\cdots[x_n\mapsto u_n].
\]
Since $u_1,\dots,u_n \in U_\Sigma$ are arbitrarily chosne, by the semantics of universal
quantifier,
\[
\mathfrak{A}^\tau_\Sigma \models \forall x_1\cdots\forall x_n\,\psi.
\]
\[
\mathfrak{A}^\tau_\Sigma \models \varphi.
\]

Since $\varphi \in \Gamma$ is also arbitrarily chosen, we have
\[
\mathfrak{A}^\tau_\Sigma \models \Gamma.
\]
\end{proof}
\bigskip

\section*{Question 4}

\begin{theorem}[Downward Löwenheim--Skolem Theorem]
Let $\Sigma$ be a countable signature and let $\Gamma$ be a set of universally quantified
FOL sentences without equality.

If $\Gamma$ is satisfiable, then there exists a \emph{countable} structure $\mathcal{A}$ such that
\[
\mathcal{A} \models \Gamma.
\]
\end{theorem}

\textbf{Question.} Prove Theorem 2.
\begin{proof}
    Suppose $\Gamma$ is satisfiable, $Ground(\Gamma)$ is propositionally sat, there exists a Herbrand structure
    \[
    \mathfrak{A}^\tau_\Sigma = (U_\Sigma, I^\tau_\Sigma)
    \]
    such that
    \[
    \mathfrak{A}^\tau_\Sigma \models \Gamma.
    \]
    Easy to show that $\mathfrak{A}^\tau_\Sigma$ is countable
    \[
    U_\Sigma = Terms_\Sigma,
    \]
    the set of all ground terms over the signature $\Sigma$. Since $\Sigma$ is countable, and every ground term is a finite string over the symbols of $\Sigma$, the set $Terms_\Sigma$ is countable.\\
    Therefore $\mathfrak{A}^\tau_\Sigma$ is a countable structure satisfying $\Gamma$.
    \end{proof}
\bigskip

\section*{Question 5}

\begin{theorem}
$\Gamma$ is satisfiable if and only if $\mathrm{Ground}(\Gamma)$ is satisfiable.
\end{theorem}

Two sets of sentences $\Gamma_1,\Gamma_2$ are said to be \emph{equivalent} if for every
structure $\mathcal{A}$,
\[
\mathcal{A}\models \Gamma_1
\iff
\mathcal{A}\models \Gamma_2 .
\]

\textbf{Question.}

Describe:

\begin{itemize}
\item a signature $\Sigma$,
\item a set $\Gamma$ of universally quantified FOL sentences without equality over $\Sigma$,
\item a $\Sigma$-structure $\mathcal{A}$,
\end{itemize}

such that

\[
\mathcal{A}\models \mathrm{Ground}(\Gamma)
\quad \text{but} \quad
\mathcal{A}\not\models \Gamma.
\]

\begin{proof}
    Consider signature
    \[\Sigma = (\{0\}, \emptyset, \{R^1\})\]
    structure 
    \[\mathfrak{A} = (\{0,1\}, I(0) = 0, I(R)=\{0\})\]
    singleton set
    \[\Gamma = \{\forall x R(x)\}\]
    \[Ground(\Gamma) = \{R(0)\}\]
    \[\mathfrak{A} \models Ground(\Gamma)\]
    but    
    \[\mathfrak{A} \not\models \Gamma\]
    as $R(1) \notin I(R)$. Basically grounding loses information about elements in the universe that no ground terms represent
\end{proof}

\bigskip

\section*{Question 6}

Let

\[
\Sigma = (\{c,d,e\}, \{f\}, \{R\})
\]

where $f$ is unary and $R$ is unary.

Define the formulas

\[
\psi_1 \equiv \neg(c=c)
\]

\[
\psi_2 \equiv (c=d) \wedge \neg(d=c)
\]

\[
\psi_3 \equiv (c=d) \wedge (d=e) \wedge \neg(c=e)
\]

\[
\psi_4 \equiv (c=d) \wedge \neg(f(c)=f(d))
\]

\[
\psi_5 \equiv (c=d) \wedge R(c) \wedge \neg R(d)
\]

Let

\[
\Gamma_i = \{\psi_i\}
\]

for $i\in\{1,2,3,4,5\}$.

\textbf{Question.}

Show that for each $i$,

\[
\mathrm{Ground}(\Gamma_i)
\]

is propositionally satisfiable but

\[
\Gamma_i
\]

is not satisfiable.
\begin{enumerate}
    \item $\psi_1 \equiv \neg(c=c)$

    \begin{itemize}
        \item Propositionally satisfiable: let $\tau(c=c)=\mathrm{false}$.
        \item Not satisfiable in FOL: equality is reflexive, $c=c$
    \end{itemize}

    \item $\psi_2 \equiv (c=d)\wedge \neg(d=c)$

    \begin{itemize}
        \item Propositionally satisfiable: let $\tau(c=d)=\mathrm{true}$ and $\tau(d=c)=\mathrm{false}$.
        \item Not satisfiable in FOL: equality is symmetric, $c=d \Rightarrow d=c$.
    \end{itemize}

    \item $\psi_3 \equiv (c=d)\wedge (d=e)\wedge \neg(c=e)$

    \begin{itemize}
        \item Propositionally satisfiable: let $\tau(c=d)=\tau(d=e)=\mathrm{true}$ and $\tau(c=e)=\mathrm{false}$.
        \item Not satisfiable in FOL: equality is transitive, $c=d \land d=e \implies c=e$.
    \end{itemize}

    \item $\psi_4 \equiv (c=d)\wedge \neg(f(c)=f(d))$

    \begin{itemize}
        \item Propositionally satisfiable: let $\tau(c=d)=\mathrm{true}$ and $\tau(f(c)=f(d))=\mathrm{false}$.
        \item Not satisfiable in FOL: functions preserve equality, $c=d \Rightarrow f(c)=f(d)$.
    \end{itemize}

    \item $\psi_5 \equiv (c=d)\wedge R(c)\wedge \neg R(d)$

    \begin{itemize}
        \item Propositionally satisfiable: let $\tau(c=d)=\mathrm{true}$, $\tau(R(c))=\mathrm{true}$, $\tau(R(d))=\mathrm{false}$.
        \item Not satisfiable in FOL: relations preserve equality, $c=d \Rightarrow (R(c)\leftrightarrow R(d))$.
    \end{itemize}
\end{enumerate}
\bigskip

\section*{Question 7}

Let

\[
\Gamma^* = eq(\Gamma) \cup E
\]

where $E$ is the set of equality axioms.

\begin{lemma}
If $\Gamma$ is satisfiable then $\Gamma^*$ is satisfiable.
\end{lemma}

\textbf{Question.} Prove Lemma 4.
\begin{proof}
    Intuition: $\Gamma \mapsto \Gamma^*$ just constraints the interpretation of relation EQ, therefore under any structure, 
    if a formula $\phi \in \Gamma$ is true, corresponding $eq(\phi) \in \Gamma^*$ is true\\
    Suppose $\Gamma$ is sat, then there exists a structure
    $\mathfrak A=(U,I)$ such that $\mathfrak A \models \Gamma$.
    
    Define a new structure $\mathfrak A'=(U,I')$ based on $\mathfrak{A}$ as such
    \begin{itemize}
        \item For all symbols in $\Sigma$, interpretation remains the same 
        \item Interpret EQ as
        \[
        I'(EQ) = \{(u,u)\mid u\in U\}.
        \]
    \end{itemize}
    
    Then $\mathfrak A'$ satisfies all equality constraints $E$, as identity is an equivalence relation and is preserved by functions and relations.
    also, since $eq(\Gamma)$ is obtained from $\Gamma$ by replacing equality with $EQ$, and $EQ$ is interpreted as identity, trivially
    \[
    \mathfrak A' \models eq(\Gamma).
    \]
    \[
    \mathfrak A' \models eq(\Gamma)\cup E = \Gamma^*.
    \]
    \end{proof}
\bigskip

\section*{Question 8}

\begin{definition}[Equivalence Relation]
A binary relation $R\subseteq U\times U$ is an \emph{equivalence relation} if it is

\begin{itemize}
\item reflexive
\item symmetric
\item transitive
\end{itemize}

For $u\in U$ the equivalence class is

\[
[u]_R = \{v\mid (u,v)\in R\}.
\]

The quotient set is

\[
U/R = \{[u]_R \mid u\in U\}.
\]

\end{definition}

\begin{definition}[Congruence Relation]
Let $U$ be a set and $E\subseteq U\times U$.

$E$ is a \emph{congruence relation} with respect to functions $F$ and relations $R$ if

\begin{enumerate}

\item $E$ is an equivalence relation.

\item For every function $f$ of arity $r$,
\[
E(u_1,v_1),\dots,E(u_r,v_r)
\Rightarrow
E(f(u_1,\dots,u_r),f(v_1,\dots,v_r)).
\]

\item For every relation $R$ of arity $r$,
\[
E(u_1,v_1),\dots,E(u_r,v_r)
\]
implies the tuples either both belong to $R$ or both do not.

\end{enumerate}

\end{definition}

\textbf{Question.}

Let $\mathcal{A}=(U,I)$ be a structure for the signature

\[
\Sigma^*=(C,F,R\cup\{EQ\})
\]

such that

\[
\mathcal{A}\models E.
\]

Show that $I(EQ)$ is a congruence relation with respect to every function symbol
$f\in F$ and relation symbol $R\in R$.
\begin{proof}
    since $\mathcal A \models E_1,E_2,E_3$, the relation $I(EQ)$ is reflexive, symmetric, and transitive, therefore $I(EQ)$ is an equivalence relation.\\
    let $f\in F$ be of arity $r$, and let $u_1,\dots,u_r,v_1,\dots,v_r \in U$ satisfy
    \[
    (u_i,v_i)\in I(EQ)\qquad (1\le i\le r).
    \]
    Since $\mathcal A \models E_4$, it follows that
    \[
    \bigl(f(u_1,\dots,u_r),\,f(v_1,\dots,v_r)\bigr)\in I(EQ).
    \]
    Thus $I(EQ)$ is preserved by every function symbol.
    
    Similarly for relation symbols we have
    \[
    R(u_1,\dots,u_r)\iff R(v_1,\dots,v_r).
    \]
    Hence $(u_1,\dots,u_r)\in I(R)$ iff $(v_1,\dots,v_r)\in I(R)$, so $I(EQ)$ is consistent with every relation symbol.
    
    Therefore $I(EQ)$ is a congruence relation with respect to every function symbol in $F$ and every relation symbol in $\mathcal R$.
    \end{proof}

\section*{Question 9}

\begin{lemma}
If $\Gamma^*$ is satisfiable then $\Gamma$ is satisfiable.
\end{lemma}

\textbf{Question.} Prove Lemma 5.
Heavy duty at work this 2 weeks, unable to finish, will attempt alongside with hw5
\bigskip

\section*{Question 10}

\begin{theorem}
Assume that the signature $\Sigma$ contains at least one constant.

Let $\Gamma$ be a set of universally quantified FOL sentences in prenex normal form.

Then

\[
\Gamma
\text{ is satisfiable }
\iff
\mathrm{Ground}(\Gamma^*)
\text{ is propositionally satisfiable.}
\]

\end{theorem}

\textbf{Question.} Prove Theorem 5.




\end{document}



